% --- Archon physics formalization source begin ---
% archon:physics
% archon:covers PhyXMiniProblems/problem_phyx_mini_0622.lean
% archon:source-report reports/phyx_mini/problem_phyx_mini_0622.source.json

\chapter{Physics problem phyx\_mini\_0622}
\label{ch:PhyXMiniProblems_problem_phyx_mini_0622}

\paragraph{Problem source.}
\#\# Physical scenario

A rectangular painting measures \textbackslash{}(1.00\textbackslash{},\textbackslash{}mathrm\{m\}\textbackslash{}) tall and \textbackslash{}(1.50\textbackslash{},\textbackslash{}mathrm\{m\}\textbackslash{}) wide, figure. It is hung on the side wall of a spaceship which is moving past the Earth at a speed of \textbackslash{}(0.90c\textbackslash{}).

\#\# Question

What are the dimensions as seen by an observer on the Earth?

\#\# Answer choices

- A: A: 0.69m

- B: B: 0.67m

- C: C: 0.65m

- D: D: 0.62m

\#\# Figure caption (auxiliary; use the image as primary evidence)

The image is a rectangular painting with a visible scene and some annotations. 

\#\#\# Content:
1. **Scene**:
   - The painting depicts a landscape with mountains and a river running through it.
   - There are trees in the foreground.
   - A sun is shown in the sky, possibly depicting either sunrise or sunset.

2. **Annotations**:
   - The height of the painting is labeled as `1.00 m`.
   - The width of the painting is indicated by a question mark, suggesting it is unknown or to be determined.

3. **Appearance and Colors**:
   - The sky is a gradient of blue, transitioning to a warmer hue near the horizon.
   - The mountains are brown, with distinctive peaks.
   - The river is a light color, likely blue, contrasting with the surrounding landscape.
   - The trees are green.
   - The sun is orange.

The painting is set against a light yellow background.

\#\# Dataset metadata

Relativity; Physical Model Grounding Reasoning, Predictive Reasoning

\paragraph{Recorded answer/context.}
C: C: 0.65m

\paragraph{Figure/image path.}
/root/proposal\_for\_physic/science-mango/phyx\_mini\_run/phyx\_data/test\_image/622.png

\paragraph{Formalization target.}
create a compiling Lean file with sorry bodies at `PhyXMiniProblems/problem\_phyx\_mini\_0622.lean`. The Lean declarations must preserve the physical quantities, dimensions or dimensional roles, figure labels, governing-law hypotheses, and final relation expressed by this problem.
Use Mathlib/Physlib names found through LeanExplore where available. If a physics API is missing, introduce faithful local abstractions rather than scalar placeholder aliases.

\begin{theorem}[Physics formalization target]
\label{thm:physics:phyx_mini_0622:target}
\lean{PhyXMiniProblems.ProblemPhyXMini0622.problem_phyx_mini_0622}
\uses{def:physics:phyx-mini-0622:phyxminiproblems-problemphyxmini0622-lengthinmeters, def:physics:phyx-mini-0622:phyxminiproblems-problemphyxmini0622-relativisticpaintingsetup, def:physics:phyx-mini-0622:phyxminiproblems-problemphyxmini0622-matchespaintingonspaceshipscenario, def:physics:phyx-mini-0622:phyxminiproblems-problemphyxmini0622-matchessuppliedpaintingfigure, def:physics:phyx-mini-0622:phyxminiproblems-problemphyxmini0622-matchesproblemreadouts, def:physics:phyx-mini-0622:phyxminiproblems-problemphyxmini0622-matchessidewallmotiongeometry, def:physics:phyx-mini-0622:phyxminiproblems-problemphyxmini0622-hasphysicalrelativisticparameters, def:physics:phyx-mini-0622:phyxminiproblems-problemphyxmini0622-satisfiesaxissensitivelorentzlengthlaw, def:physics:phyx-mini-0622:phyxminiproblems-problemphyxmini0622-matchesdisplayedwidthchoice}
The assigned autoformalize agent should translate this physics problem into Lean declarations in the covered file, with theorem and lemma proof bodies written as `by sorry`.
\end{theorem}
\begin{proof}
This is an autoformalization task, not a proof task. Produce statements that can later be proved by the physics prover without weakening the physical model.
\end{proof}

% --- Archon declaration topology begin ---
\section{Declaration topology}
\begin{definition}[Length Quantity]
  \label{def:physics:phyx-mini-0622:phyxminiproblems-problemphyxmini0622-lengthquantity}
  \lean{PhyXMiniProblems.ProblemPhyXMini0622.LengthQuantity}
  A nonnegative, unit-independent physical length
\end{definition}
\begin{proof}
  This is the declared model definition of Length Quantity.
\end{proof}
\begin{definition}[length Readout]
  \label{def:physics:phyx-mini-0622:phyxminiproblems-problemphyxmini0622-lengthreadout}
  \lean{PhyXMiniProblems.ProblemPhyXMini0622.lengthReadout}
  \uses{def:physics:phyx-mini-0622:phyxminiproblems-problemphyxmini0622-lengthquantity}
  Read a physical length in a selected length unit
\end{definition}
\begin{proof}
  This is the declared model definition of length Readout.
\end{proof}
\begin{definition}[length In Meters]
  \label{def:physics:phyx-mini-0622:phyxminiproblems-problemphyxmini0622-lengthinmeters}
  \lean{PhyXMiniProblems.ProblemPhyXMini0622.lengthInMeters}
  \uses{def:physics:phyx-mini-0622:phyxminiproblems-problemphyxmini0622-lengthquantity, def:physics:phyx-mini-0622:phyxminiproblems-problemphyxmini0622-lengthreadout}
  Meter readout used by the problem, figure, and answer choices
\end{definition}
\begin{proof}
  This is the declared model definition of length In Meters.
\end{proof}
\begin{definition}[speed In Meters Per Second]
  \label{def:physics:phyx-mini-0622:phyxminiproblems-problemphyxmini0622-speedinmeterspersecond}
  \lean{PhyXMiniProblems.ProblemPhyXMini0622.speedInMetersPerSecond}
  Read a physical speed in metres per second
\end{definition}
\begin{proof}
  This is the declared model definition of speed In Meters Per Second.
\end{proof}
\begin{definition}[vacuum Speed Of Light In Meters Per Second]
  \label{def:physics:phyx-mini-0622:phyxminiproblems-problemphyxmini0622-vacuumspeedoflightinmeterspersecond}
  \lean{PhyXMiniProblems.ProblemPhyXMini0622.vacuumSpeedOfLightInMetersPerSecond}
  Physlib's vacuum speed of light, read in metres per second
\end{definition}
\begin{proof}
  This is the declared model definition of vacuum Speed Of Light In Meters Per Second.
\end{proof}
\begin{definition}[Inertial Frame Label]
  \label{def:physics:phyx-mini-0622:phyxminiproblems-problemphyxmini0622-inertialframelabel}
  \lean{PhyXMiniProblems.ProblemPhyXMini0622.InertialFrameLabel}
  The two inertial frames relevant to the dimension measurements
\end{definition}
\begin{proof}
  This is the declared model definition of Inertial Frame Label.
\end{proof}
\begin{definition}[Painting Dimension]
  \label{def:physics:phyx-mini-0622:phyxminiproblems-problemphyxmini0622-paintingdimension}
  \lean{PhyXMiniProblems.ProblemPhyXMini0622.PaintingDimension}
  The two physical dimensions of the rectangular painting
\end{definition}
\begin{proof}
  This is the declared model definition of Painting Dimension.
\end{proof}
\begin{definition}[Motion Axis Role]
  \label{def:physics:phyx-mini-0622:phyxminiproblems-problemphyxmini0622-motionaxisrole}
  \lean{PhyXMiniProblems.ProblemPhyXMini0622.MotionAxisRole}
  Whether a painting dimension is parallel or transverse to the motion
\end{definition}
\begin{proof}
  This is the declared model definition of Motion Axis Role.
\end{proof}
\begin{definition}[Figure Dimension Annotation]
  \label{def:physics:phyx-mini-0622:phyxminiproblems-problemphyxmini0622-figuredimensionannotation}
  \lean{PhyXMiniProblems.ProblemPhyXMini0622.FigureDimensionAnnotation}
  The two kinds of dimension annotation appearing in the supplied image
\end{definition}
\begin{proof}
  This is the declared model definition of Figure Dimension Annotation.
\end{proof}
\begin{definition}[Painting Figure]
  \label{def:physics:phyx-mini-0622:phyxminiproblems-problemphyxmini0622-paintingfigure}
  \lean{PhyXMiniProblems.ProblemPhyXMini0622.PaintingFigure}
  \uses{def:physics:phyx-mini-0622:phyxminiproblems-problemphyxmini0622-paintingdimension, def:physics:phyx-mini-0622:phyxminiproblems-problemphyxmini0622-figuredimensionannotation}
  Literal evidence carried by phyx\_data/test\_image/622.png. The image labels the vertical arrow as 1.00 m and places a question mark under the horizontal arrow. The landscape content is retained as a qualitative check that the outlined rectangle is the painting referred to in the prose
\end{definition}
\begin{proof}
  This is the declared model definition of Painting Figure.
\end{proof}
\begin{definition}[Relativistic Painting Setup]
  \label{def:physics:phyx-mini-0622:phyxminiproblems-problemphyxmini0622-relativisticpaintingsetup}
  \lean{PhyXMiniProblems.ProblemPhyXMini0622.RelativisticPaintingSetup}
  \uses{def:physics:phyx-mini-0622:phyxminiproblems-problemphyxmini0622-lengthquantity, def:physics:phyx-mini-0622:phyxminiproblems-problemphyxmini0622-inertialframelabel, def:physics:phyx-mini-0622:phyxminiproblems-problemphyxmini0622-paintingdimension, def:physics:phyx-mini-0622:phyxminiproblems-problemphyxmini0622-motionaxisrole, def:physics:phyx-mini-0622:phyxminiproblems-problemphyxmini0622-paintingfigure}
  Independent physical quantities in the problem. In particular, earthMeasuredDimension is an observable field rather than a definition made from an answer choice or from the contraction formula
\end{definition}
\begin{proof}
  This is the declared model definition of Relativistic Painting Setup.
\end{proof}
\begin{definition}[speed Fraction Of Light]
  \label{def:physics:phyx-mini-0622:phyxminiproblems-problemphyxmini0622-speedfractionoflight}
  \lean{PhyXMiniProblems.ProblemPhyXMini0622.speedFractionOfLight}
  \uses{def:physics:phyx-mini-0622:phyxminiproblems-problemphyxmini0622-speedinmeterspersecond, def:physics:phyx-mini-0622:phyxminiproblems-problemphyxmini0622-vacuumspeedoflightinmeterspersecond, def:physics:phyx-mini-0622:phyxminiproblems-problemphyxmini0622-relativisticpaintingsetup}
  The dimensionless speed parameter β = v/c in coherent SI readouts
\end{definition}
\begin{proof}
  This is the declared model definition of speed Fraction Of Light.
\end{proof}
\begin{definition}[lorentz Factor]
  \label{def:physics:phyx-mini-0622:phyxminiproblems-problemphyxmini0622-lorentzfactor}
  \lean{PhyXMiniProblems.ProblemPhyXMini0622.lorentzFactor}
  \uses{def:physics:phyx-mini-0622:phyxminiproblems-problemphyxmini0622-relativisticpaintingsetup, def:physics:phyx-mini-0622:phyxminiproblems-problemphyxmini0622-speedfractionoflight}
  The Lorentz factor supplied by Physlib for the setup's relative speed
\end{definition}
\begin{proof}
  This is the declared model definition of lorentz Factor.
\end{proof}
\begin{definition}[Matches Painting On Spaceship Scenario]
  \label{def:physics:phyx-mini-0622:phyxminiproblems-problemphyxmini0622-matchespaintingonspaceshipscenario}
  \lean{PhyXMiniProblems.ProblemPhyXMini0622.MatchesPaintingOnSpaceshipScenario}
  \uses{def:physics:phyx-mini-0622:phyxminiproblems-problemphyxmini0622-relativisticpaintingsetup}
  The prose scenario and the operational meaning of an Earth-frame length
\end{definition}
\begin{proof}
  This is the declared model definition of Matches Painting On Spaceship Scenario.
\end{proof}
\begin{definition}[Matches Supplied Painting Figure]
  \label{def:physics:phyx-mini-0622:phyxminiproblems-problemphyxmini0622-matchessuppliedpaintingfigure}
  \lean{PhyXMiniProblems.ProblemPhyXMini0622.MatchesSuppliedPaintingFigure}
  \uses{def:physics:phyx-mini-0622:phyxminiproblems-problemphyxmini0622-paintingfigure}
  Quantitative and qualitative evidence read directly from the bitmap
\end{definition}
\begin{proof}
  This is the declared model definition of Matches Supplied Painting Figure.
\end{proof}
\begin{definition}[Matches Problem Readouts]
  \label{def:physics:phyx-mini-0622:phyxminiproblems-problemphyxmini0622-matchesproblemreadouts}
  \lean{PhyXMiniProblems.ProblemPhyXMini0622.MatchesProblemReadouts}
  \uses{def:physics:phyx-mini-0622:phyxminiproblems-problemphyxmini0622-lengthinmeters, def:physics:phyx-mini-0622:phyxminiproblems-problemphyxmini0622-relativisticpaintingsetup, def:physics:phyx-mini-0622:phyxminiproblems-problemphyxmini0622-speedfractionoflight}
  Numerical data stated in the problem. The height is linked to the primary figure's printed rest-frame value; the width and speed come from the prose. No Earth-frame dimension or answer-choice value occurs here
\end{definition}
\begin{proof}
  This is the declared model definition of Matches Problem Readouts.
\end{proof}
\begin{definition}[Matches Side Wall Motion Geometry]
  \label{def:physics:phyx-mini-0622:phyxminiproblems-problemphyxmini0622-matchessidewallmotiongeometry}
  \lean{PhyXMiniProblems.ProblemPhyXMini0622.MatchesSideWallMotionGeometry}
  \uses{def:physics:phyx-mini-0622:phyxminiproblems-problemphyxmini0622-relativisticpaintingsetup}
  The side-wall orientation relevant to Lorentz length contraction
\end{definition}
\begin{proof}
  This is the declared model definition of Matches Side Wall Motion Geometry.
\end{proof}
\begin{definition}[Has Physical Relativistic Parameters]
  \label{def:physics:phyx-mini-0622:phyxminiproblems-problemphyxmini0622-hasphysicalrelativisticparameters}
  \lean{PhyXMiniProblems.ProblemPhyXMini0622.HasPhysicalRelativisticParameters}
  \uses{def:physics:phyx-mini-0622:phyxminiproblems-problemphyxmini0622-lengthinmeters, def:physics:phyx-mini-0622:phyxminiproblems-problemphyxmini0622-vacuumspeedoflightinmeterspersecond, def:physics:phyx-mini-0622:phyxminiproblems-problemphyxmini0622-relativisticpaintingsetup, def:physics:phyx-mini-0622:phyxminiproblems-problemphyxmini0622-speedfractionoflight}
  Positivity and the nonnegative subluminal regime of the physical model
\end{definition}
\begin{proof}
  This is the declared model definition of Has Physical Relativistic Parameters.
\end{proof}
\begin{definition}[Satisfies Axis Sensitive Lorentz Length Law]
  \label{def:physics:phyx-mini-0622:phyxminiproblems-problemphyxmini0622-satisfiesaxissensitivelorentzlengthlaw}
  \lean{PhyXMiniProblems.ProblemPhyXMini0622.SatisfiesAxisSensitiveLorentzLengthLaw}
  \uses{def:physics:phyx-mini-0622:phyxminiproblems-problemphyxmini0622-lengthreadout, def:physics:phyx-mini-0622:phyxminiproblems-problemphyxmini0622-paintingdimension, def:physics:phyx-mini-0622:phyxminiproblems-problemphyxmini0622-relativisticpaintingsetup, def:physics:phyx-mini-0622:phyxminiproblems-problemphyxmini0622-lorentzfactor}
  The governing special-relativistic length law. For every painting dimension and every common length unit, a component parallel to the relative motion is divided by γ(β), while a transverse component is unchanged. This generic law contains neither the numerical Earth-frame dimensions nor an answer label
\end{definition}
\begin{proof}
  This is the declared model definition of Satisfies Axis Sensitive Lorentz Length Law.
\end{proof}
\begin{definition}[Answer Choice]
  \label{def:physics:phyx-mini-0622:phyxminiproblems-problemphyxmini0622-answerchoice}
  \lean{PhyXMiniProblems.ProblemPhyXMini0622.AnswerChoice}
  Labels of the four width choices printed with the problem
\end{definition}
\begin{proof}
  This is the declared model definition of Answer Choice.
\end{proof}
\begin{definition}[displayed Width In Meters]
  \label{def:physics:phyx-mini-0622:phyxminiproblems-problemphyxmini0622-displayedwidthinmeters}
  \lean{PhyXMiniProblems.ProblemPhyXMini0622.displayedWidthInMeters}
  \uses{def:physics:phyx-mini-0622:phyxminiproblems-problemphyxmini0622-answerchoice}
  Width in metres printed beside each answer label
\end{definition}
\begin{proof}
  This is the declared model definition of displayed Width In Meters.
\end{proof}
\begin{definition}[recorded Dataset Answer]
  \label{def:physics:phyx-mini-0622:phyxminiproblems-problemphyxmini0622-recordeddatasetanswer}
  \lean{PhyXMiniProblems.ProblemPhyXMini0622.recordedDatasetAnswer}
  \uses{def:physics:phyx-mini-0622:phyxminiproblems-problemphyxmini0622-answerchoice}
  The answer label recorded by the source dataset; this is metadata only
\end{definition}
\begin{proof}
  This is the declared model definition of recorded Dataset Answer.
\end{proof}
\begin{definition}[Matches Displayed Width Choice]
  \label{def:physics:phyx-mini-0622:phyxminiproblems-problemphyxmini0622-matchesdisplayedwidthchoice}
  \lean{PhyXMiniProblems.ProblemPhyXMini0622.MatchesDisplayedWidthChoice}
  \uses{def:physics:phyx-mini-0622:phyxminiproblems-problemphyxmini0622-lengthinmeters, def:physics:phyx-mini-0622:phyxminiproblems-problemphyxmini0622-relativisticpaintingsetup, def:physics:phyx-mini-0622:phyxminiproblems-problemphyxmini0622-answerchoice, def:physics:phyx-mini-0622:phyxminiproblems-problemphyxmini0622-displayedwidthinmeters}
  Agreement with a width displayed to two decimal places. The strict tolerance 0.005 m excludes a rounding tie at a hundredth-place boundary
\end{definition}
\begin{proof}
  This is the declared model definition of Matches Displayed Width Choice.
\end{proof}
\begin{lemma}[earth Measured Height In Meters eq one]
  \label{lem:physics:phyx-mini-0622:phyxminiproblems-problemphyxmini0622-earthmeasuredheightinmeters-eq-one}
  \lean{PhyXMiniProblems.ProblemPhyXMini0622.earthMeasuredHeightInMeters_eq_one}
  \uses{def:physics:phyx-mini-0622:phyxminiproblems-problemphyxmini0622-lengthinmeters, def:physics:phyx-mini-0622:phyxminiproblems-problemphyxmini0622-relativisticpaintingsetup, def:physics:phyx-mini-0622:phyxminiproblems-problemphyxmini0622-matchessuppliedpaintingfigure, def:physics:phyx-mini-0622:phyxminiproblems-problemphyxmini0622-matchesproblemreadouts, def:physics:phyx-mini-0622:phyxminiproblems-problemphyxmini0622-matchessidewallmotiongeometry, def:physics:phyx-mini-0622:phyxminiproblems-problemphyxmini0622-satisfiesaxissensitivelorentzlengthlaw}
  The transverse 1.00 m height is unchanged for the Earth observer
\end{lemma}
\begin{proof}
  The conclusion follows from the governing relations and figure readouts named in the statement of earth Measured Height In Meters eq one.
\end{proof}
\begin{lemma}[earth Measured Width In Meters exact]
  \label{lem:physics:phyx-mini-0622:phyxminiproblems-problemphyxmini0622-earthmeasuredwidthinmeters-exact}
  \lean{PhyXMiniProblems.ProblemPhyXMini0622.earthMeasuredWidthInMeters_exact}
  \uses{def:physics:phyx-mini-0622:phyxminiproblems-problemphyxmini0622-lengthinmeters, def:physics:phyx-mini-0622:phyxminiproblems-problemphyxmini0622-relativisticpaintingsetup, def:physics:phyx-mini-0622:phyxminiproblems-problemphyxmini0622-matchesproblemreadouts, def:physics:phyx-mini-0622:phyxminiproblems-problemphyxmini0622-matchessidewallmotiongeometry, def:physics:phyx-mini-0622:phyxminiproblems-problemphyxmini0622-satisfiesaxissensitivelorentzlengthlaw}
  The longitudinal Earth-frame width is the 1.50 m proper width divided by the Lorentz factor at β = 0.90
\end{lemma}
\begin{proof}
  The conclusion follows from the governing relations and figure readouts named in the statement of earth Measured Width In Meters exact.
\end{proof}
% --- Archon declaration topology end ---
% --- Archon physics formalization source end ---
